\chapter{Выводы}
\label{chap:conclusion}

Гипотеза о влиянии глубины подтвердилась лишь частично: ёмкость сети управляет поэлементной точностью, но структурное качество от неё почти не зависит и при чисто пиксельном обучении остаётся ниже приемлемого. Гипотеза о состязательном дообучении верна без оговорок, причём свойственный классическим SR-GAN компромисс между точностью и резкостью здесь не возникает: добавление физического регуляризатора удерживает обе группы метрик. Отсюда главный практический вывод: самая компактная модель не уступает тяжёлым, а по структуре даже превосходит их.

Достигнутый уровень структурного качества, по-видимому, упирается в небольшой обучающий набор, поэтому первоочередным направлением дальнейшей работы остаётся его расширение и пополнение более разнообразными режимами фильтрации.